\documentclass[11pt]{article}

\usepackage[utf8]{inputenc}
\usepackage[T1]{fontenc}
\usepackage[margin=1in]{geometry}
\usepackage{amsmath,amssymb}
\usepackage{booktabs}
\usepackage{graphicx}
\usepackage{natbib}
\usepackage[hidelinks]{hyperref}
% microtype evens out inter-word spacing in justified text (character
% protrusion + font expansion) — the standard fix for uneven gaps.
\usepackage{microtype}

\title{\textbf{Behavioral Cloning for Self-Driving Cars:\\
A Leak-Free Re-evaluation and a Negative Result on Temporal Models}}

\author{Aly Elgemei\\
  Independent Researcher, Siegen, Germany\\
  \url{https://github.com/ali7996/behavioral-cloning-thesis}}

\date{2026}

\begin{document}
\maketitle

\begin{abstract}
Behavioral cloning trains a neural network to imitate a human driver by mapping a
forward-facing camera image to a steering command. We revisit the practical
component of a 2023 Master's thesis on this task~\citep{elgemei2023behavioral} and
report a corrected, honestly-evaluated reproduction. Three issues are identified and
addressed. First, the thesis's two recurrent models were inadvertently trained on
uniform random noise rather than on driving data; we retrain them on real ten-frame
image sequences. Second, the thesis's 2D-CNN+LSTM consumed a single frame reshaped
into a length-one sequence and was therefore not a temporal model at all; we rebuild
it as a genuine per-frame convolutional encoder followed by a recurrent head. Third,
and most consequential, the conventional random train/validation split leaks
information: consecutive frames are roughly 70\,ms apart and nearly identical, so a
random split places near-duplicates of validation frames into the training set.
Re-evaluating the convolutional baseline on a leak-free contiguous split lowers its
coefficient of determination from $0.52$ to $0.31$. On that leak-free split, neither
temporal model beats the single-frame convolutional baseline; in fact neither beats
predicting the mean steering angle ($R^2 \le 0$). We report this negative result
as-is: on this dataset, at this scale, temporal context does not improve steering
prediction. All code, trained weights, and evaluation scripts are public.
\end{abstract}

\section{Introduction}

End-to-end behavioral cloning is the oldest and simplest paradigm for vision-based
autonomous driving: a neural network is trained, by supervised regression, to
reproduce the control output of a human demonstrator directly from raw camera
input~\citep{pomerleau1989alvinn,bojarski2016end}. It requires no hand-engineered
perception stack and no explicit map, which makes it an attractive teaching and
benchmarking task --- and a frequent subject of student projects and theses.

This paper revisits one such project: the practical component of a 2023 Master's
thesis~\citep{elgemei2023behavioral} that implemented and compared three
architectures for predicting the steering angle of a car in the Udacity driving
simulator. While re-examining the original code in order to publish it, we found
methodological flaws that invalidated part of the thesis's reported results. Rather
than quietly discard the affected numbers, we treat the re-examination itself as the
contribution: we correct the flaws, re-run the experiments under a leak-free
evaluation protocol, and report the outcome --- including a negative result ---
transparently.

\paragraph{Contributions.}
\begin{enumerate}
  \item We reproduce the convolutional baseline on the real driving dataset and
        find that, retrained honestly, it roughly halves the steering error
        reported in the thesis.
  \item We document that the thesis's two recurrent models were trained on
        \texttt{numpy} uniform-random arrays rather than on driving data, so their
        reported metrics measured nothing, and we retrain them on real ten-frame
        sequences.
  \item We show that the conventional random frame split leaks information between
        training and validation, because adjacent frames in a continuous drive are
        near-duplicates; we introduce a contiguous, timeline-ordered split that
        removes the leak, and quantify its effect.
  \item Under the leak-free protocol we report a negative result: temporal
        (recurrent) models do not improve on the single-frame convolutional
        baseline, and do not beat a constant mean predictor.
\end{enumerate}

The intent of this report is not to claim a new architecture or a new
state-of-the-art. It is to demonstrate that careful, honest evaluation of a small
and familiar task can still produce a useful and non-obvious finding.

\section{Background and Related Work}

\paragraph{End-to-end driving.}
The idea of learning to steer directly from pixels dates to ALVINN, which trained a
small fully-connected network to drive a vehicle from camera and range
input~\citep{pomerleau1989alvinn}. The approach was revived at scale by NVIDIA's
DAVE-2 system, in which a convolutional network maps a single front-facing camera
frame to a steering command and learns useful road features without explicit
supervision for them~\citep{bojarski2016end}. The architecture used in this paper is
the DAVE-2 network, and the dataset is generated in a driving simulator of the same
family.

\paragraph{Behavioral cloning and its limits.}
Plain behavioral cloning is known to be brittle. \citet{codevilla2018end} show that
an unconditioned policy cannot resolve the ambiguity of an intersection --- the same
image admits a left or a right turn --- and condition the policy on a navigational
command to fix this. \citet{codevilla2019exploring} further document that behavior
cloning suffers from high variance across training seeds, sensitivity to dataset
composition, and a driving-performance ceiling that more data does not lift. Our
negative result on temporal models is consistent with this literature: behavior
cloning's difficulties are not obviously solved by adding model capacity or input
context.

\paragraph{Recurrent models for sequences.}
Long Short-Term Memory networks~\citep{hochreiter1997long} are the standard tool for
modelling temporal dependencies, and it is natural to ask whether a short window of
consecutive frames lets a driving model exploit motion that a single frame cannot.
The thesis posed exactly this question. We show below that answering it requires
both real sequential data and a split that does not leak temporally-adjacent frames
--- neither of which the original implementation had.

\paragraph{Leakage in sequential data.}
When samples are drawn from a continuous process, examples that are close in time
are statistically dependent. A uniformly random train/validation split then places
near-identical examples on both sides of the partition, and the validation score
measures memorization of almost-seen inputs rather than generalization. This is a
well-understood pitfall for time-series and video data; we quantify its size for the
present task in Section~\ref{sec:leakage}.

\section{Task, Dataset, and Preprocessing}

\paragraph{Task.}
Given a forward-facing camera image (or a short window of consecutive images), the
model predicts a scalar steering angle in $[-1, 1]$. This is a regression problem,
and we evaluate it with regression metrics. We deliberately drop the pseudo-accuracy
$1/(1+\mathrm{MSE})$ used in the original thesis: it is a monotone transform of the
error, not a meaningful ``accuracy'', and it makes weak models look strong.

\paragraph{Dataset.}
We use the public Udacity simulator dataset \texttt{rslim087a/track}: 4{,}053 frames
recorded on a single circuit, with center, left, and right camera images and
synchronized telemetry. Inspecting the timestamps encoded in the frame filenames, we
find the recording is one uninterrupted run at roughly 14\,Hz --- the largest gap
between consecutive frames is 0.48\,s, and there are no gaps above 0.5\,s. This
single-run structure is what makes the leakage analysis of Section~\ref{sec:leakage}
both necessary and clean.

\paragraph{Steering balancing.}
The raw steering distribution is dominated by near-zero (straight-ahead) values. We
bin the steering range into 25 bins and cap each bin at 400 samples, which reduces
the dataset to 1{,}463 frames and prevents the model from collapsing onto the
constant zero prediction.

\paragraph{Side cameras.}
Each frame carries three cameras. We use all three: the left and right images are
added with a fixed steering correction of $\pm 0.15$, which both triples the data
(to 4{,}389 image/steering pairs) and teaches recovery toward the lane center.

\paragraph{Preprocessing.}
Every image is cropped to remove sky and car hood, converted from RGB to YUV,
blurred with a $3\times3$ Gaussian kernel, resized to $200\times66$, and scaled to
$[0,1]$. During training only, four augmentations are each applied independently
with probability $0.5$: random zoom, pan, brightness change, and horizontal flip
(with a sign flip of the steering label).

\section{Models}
\label{sec:models}

We study three architectures. All are trained with the Adam
optimizer~\citep{kingma2015adam} ($\text{lr}=10^{-3}$) and a mean-squared-error
loss.

\paragraph{Convolutional baseline (single frame).}
This is the DAVE-2 network~\citep{bojarski2016end}: five convolutional layers
(filter counts $24$--$36$--$48$--$64$--$64$) followed by three fully-connected
layers and a scalar output, with ELU activations and dropout. It has $252{,}219$
trainable parameters and consumes one preprocessed $66\times200\times3$ frame.

\paragraph{Flatten-LSTM (ten frames).}
This is the thesis's recurrent design: a window of ten consecutive frames, each
flattened to a $39{,}600$-dimensional vector, processed by two stacked
LSTM~\citep{hochreiter1997long} layers of 128 units, then a dense head. Flattening
the image discards all spatial structure and makes the first recurrent layer
enormous; the model has $20{,}506{,}113$ parameters. We retain this design as-is so
that the temporal experiment is a faithful correction of the thesis rather than a
redesign.

\paragraph{CNN-LSTM (ten frames).}
The thesis's third model was described as a 2D-CNN+LSTM, but its implementation
applied the convolutional stack to a \emph{single} frame and then reshaped the
result into a sequence of length one before the LSTM --- so the recurrent layers saw
no temporal axis and the model was not temporal at all. We rebuild it correctly: a
strided convolutional encoder is applied per frame (time-distributed) across the
ten-frame window, producing a genuine length-ten feature sequence, which an LSTM
then summarizes. The corrected model has $1{,}105{,}701$ parameters.

\section{Two Corrections to the Evaluation}
\label{sec:corrections}

\subsection{Training on real sequences}

In the original code, the two recurrent models were trained not on driving data but
on an array of uniform random noise --- in effect, a call to
\texttt{numpy.random.random} producing a tensor of shape
$4389 \times 10 \times 66 \times 200 \times 3$ --- paired with the real
steering labels. A network trained on random inputs cannot learn an input--output
mapping, so the LSTM and 2D-CNN+LSTM metrics reported in the thesis do not measure
learning from driving data. Notably, those metrics were \emph{better} than the
convolutional baseline's, which alone should have been a warning sign.

We correct this by building real temporal windows. Frames are sorted by their
timestamp, split into maximal runs with no inter-frame gap above $0.5$\,s (here, a
single run), and a window of ten consecutive frames is slid along each run. Each
window is labelled with the steering angle of its final frame. Windows are built
independently for the three cameras, so a window is always single-camera and remains
a coherent short video. Augmentation is drawn once per window and applied identically
to all ten frames, so temporal coherence is preserved.

\subsection{A leak-free contiguous split}
\label{sec:leakage}

At 14\,Hz, consecutive frames are about 70\,ms apart and visually almost identical.
A uniformly random train/validation split therefore places, for almost every
validation frame, one or two near-duplicate frames into the training set. The model
can score well on validation by recognizing inputs it has effectively already seen.

This matters for any of the three models, but it is fatal for the temporal
experiment: ten-frame windows sampled one step apart overlap in nine of ten frames,
so a random split over windows leaks almost completely.

We therefore order all frames by time and split the \emph{timeline}: the first 80\%
of the run becomes the training region and the last 20\% the validation region.
Windows are built within each region only, so no window crosses the boundary and no
validation window shares a frame with a training window. For the temporal models we
balance the training windows by steering bin \emph{after} windowing (balancing
before windowing would punch holes in the timeline). For a fair comparison, the
convolutional baseline is then re-scored on this same contiguous validation set, by
predicting on the final frame of every validation window.

\section{Experiments}

\paragraph{Setup.}
All models are trained on a single CPU. The convolutional baseline is trained for 10
epochs. The temporal models are trained with early stopping on the validation loss;
both stopped within a few epochs. Reported metrics are mean squared error (MSE),
mean absolute error (MAE), root mean squared error (RMSE), and the coefficient of
determination $R^2$. An $R^2$ of $0$ corresponds to predicting the mean steering
angle of the validation set; a negative $R^2$ is worse than that constant baseline.

\subsection{Convolutional baseline vs.\ the thesis}

Trained on the real data with the original preprocessing and augmentation intact,
and evaluated under the thesis's own protocol (a random 80/20 split), the
convolutional baseline roughly halves the error the thesis reported
(Table~\ref{tab:repro}).

\begin{table}[h]
\centering
\caption{Convolutional baseline under the thesis's own random-split protocol.}
\label{tab:repro}
\begin{tabular}{lrrrr}
\toprule
 & MSE & MAE & RMSE & $R^2$ \\
\midrule
This work (real data, random split) & 0.054 & \textbf{0.168} & \textbf{0.233} & 0.520 \\
Thesis~\citep{elgemei2023behavioral}, reported & --- & 0.300 & 0.408 & --- \\
\bottomrule
\end{tabular}
\end{table}

\subsection{Leak-free comparison of all three models}

Table~\ref{tab:main} reports all three models on the identical leak-free contiguous
validation set ($2{,}406$ windows). The result is unambiguous: the single-frame
convolutional baseline is the only model that explains any steering variance at all.
Both temporal models have a negative $R^2$ --- they are beaten by predicting the
mean steering angle --- despite one of them having eighty times the baseline's
parameters.

\begin{table}[h]
\centering
\caption{All three models on the same leak-free contiguous validation set
($2{,}406$ windows). Lower is better for MSE/MAE/RMSE; higher is better for $R^2$.}
\label{tab:main}
\begin{tabular}{lrrrrr}
\toprule
Model & Parameters & MSE & MAE & RMSE & $R^2$ \\
\midrule
Convolutional baseline (1 frame)  & 252{,}219      & \textbf{0.019} & \textbf{0.106} & \textbf{0.138} & $\mathbf{+0.308}$ \\
Flatten-LSTM (10 frames)          & 20{,}506{,}113 & 0.029          & 0.128          & 0.169          & $-0.030$ \\
CNN-LSTM (10 frames)              & 1{,}105{,}701  & 0.029          & 0.138          & 0.170          & $-0.047$ \\
\bottomrule
\end{tabular}
\end{table}

\subsection{The size of the leakage effect}

Table~\ref{tab:leak} isolates the leakage by evaluating the \emph{same} trained
convolutional model under the two splits. Moving from the random split to the
leak-free contiguous split lowers $R^2$ from $0.52$ to $0.31$: about two-fifths of
the apparent explained variance under the random split was an artifact of
near-duplicate frames crossing the partition. (The MAE figures are not comparable
across the two rows, because the random and contiguous validation sets have
different steering distributions; $R^2$, being normalized by each set's own
variance, is the meaningful cross-split quantity.)

\begin{table}[h]
\centering
\caption{The same convolutional model evaluated under two splits. The random split
inflates the apparent fit.}
\label{tab:leak}
\begin{tabular}{lrrr}
\toprule
Split & Validation frames & RMSE & $R^2$ \\
\midrule
Random 80/20            & 800   & 0.233 & 0.520 \\
Contiguous (leak-free)  & 2{,}406 & 0.138 & 0.308 \\
\bottomrule
\end{tabular}
\end{table}

\section{Discussion}

\paragraph{Why the temporal models fail.}
The failure is not a bug in the corrected pipeline; it is a genuine property of the
task at this scale. After balancing, only about $2{,}300$ ten-frame windows remain,
and because windows are slid one step at a time they overlap heavily --- the number
of effectively independent short clips is far smaller still. Against this, the
flatten-LSTM carries twenty million parameters. Its training loss never fell below
roughly $0.18$ MSE: the model did not even fit the training set well, so it is
underfitting, not merely overfitting. The plausible explanation is that on this
single, fixed circuit the correct steering angle is almost fully determined by the
appearance of the current frame; a short window of recent frames adds little that
the current frame does not already contain, and the recurrent models pay for the
extra input dimensionality without a corresponding signal to learn from.

\paragraph{A negative result is still a result.}
The honest conclusion is that, for this dataset and this scale, temporal context
does not help, and the simple convolutional baseline should be preferred. This
contradicts the original thesis, whose recurrent models appeared to win --- but only
because they were evaluated on noise. Reporting the corrected outcome, rather than
the flattering one, is the point of this paper.

\paragraph{Random splits flatter sequential models.}
The leakage measurement is a transferable warning. Any project that draws training
examples from a continuous recording --- driving logs, video, sensor streams --- and
then splits them uniformly at random will over-report generalization. A
timeline-ordered split costs nothing and removes the illusion.

\paragraph{Limitations.}
The study uses a single simulated circuit, one human demonstrator, and a few
thousand frames. A larger and more varied corpus --- multiple tracks, real driving,
harder maneuvers --- could change the balance: temporal context may well help when
there is enough data to learn it and enough situational variety to require it. The
negative result is a statement about this regime, not a universal claim. Evaluation
is also offline; closed-loop driving in the simulator is supported by the released
code but is not scored here.

\section{Conclusion}

We revisited a Master's-thesis behavioral-cloning project, found that its recurrent
models had been trained on random noise and that its evaluation split leaked
temporally-adjacent frames, and corrected both. Retrained honestly, the convolutional
baseline roughly halves the thesis's reported error; evaluated without leakage, it is
the only one of the three models that explains any steering variance, while the two
temporal models fail to beat a constant predictor. The contribution is not a new
method but a demonstration that disciplined, leak-free evaluation of a small and
familiar task yields a clear, reportable, and initially counter-intuitive finding.
All code, trained weights, and evaluation scripts are available at
\url{https://github.com/ali7996/behavioral-cloning-thesis}.

\paragraph{Reproducibility.}
The repository contains the data pipeline, all three model definitions, training and
evaluation scripts, the contiguous-split sequence builder, and the saved metric
files from which every number in this paper is taken. The dataset is the public
\texttt{rslim087a/track} collection.

\bibliographystyle{plainnat}
\bibliography{refs}

\end{document}
